\section{证据必须控制大语言模型的代码编辑}

大语言模型（LLM）智能体可以编辑不熟悉的高层次综合（HLS）内核，但只有
工具报告能够证明正确性和硬件质量。因此，一个有计量能力的智能体需要一套
提升规则：当编辑后的代码未能通过仿真、综合、时序、容量或协同仿真中任一
必需验证门时，保留已验证的备选方案~\cite{amd2025vitis}。

现有 HLS 智能体实现了代码变换、指令搜索和反馈驱动的修复自动化
~\cite{xiong2024hlspilot,xu2024ralad,li2026chathls}。
基于工具的语言智能体也交替进行推理和行动
~\cite{yao2023react}。我们的设计将候选方案的提升决策交由
确定性的 HLS 证据控制，并在明确的模型和工具预算约束下运行。

我们的已验证候选回路（\VCL）贡献了三项构件和发现。
\textbf{（1）验证契约。} 一个 fail-closed 工作流在七个门限前保留已验证备选方案：
接口（Interface）、C 仿真（\CSim）、HLS 综合（\Synth）、频率（Frequency）、
容量（Capacity）、所需的 C/RTL 协同仿真（\CoSim）和指标完备性（Metric completeness）。
\textbf{（2）硬件目标。} 有界分数 \QHW{} 拒绝那些在时钟或最差资源质量上的损失
超过延迟收益的加速方案。
\textbf{（3）端点评估。} 第~\ref{sec:evaluation}~节将任务完成与评分覆盖分离，
对比三个云端端点。DeepSeek 在完成覆盖上领先，Qwen3.5 拥有最大的可评分集合和
最高总分，Qwen3.6 在结构修复上仍存在差距。评分覆盖不均和残余的 API 失败
阻碍了将端点差异完全归因于模型权重。
